% Part: counterfactuals
% Chapter: introduction
% Section: material-conditional

\documentclass[../../../include/open-logic-section]{subfiles}

\begin{document}

\olfileid{cnt}{int}{mat}

\olsection{భౌతిక సోపాధికం}

ఆంగ్లంలో సోపాధిక వాక్యం అత్యంత సరళమైన రూపంలో
``ఒకవేళ \dots అయితే \dots'' అనే రూపంలో ఉంటుంది;
ఆ ఖాళీల్లో కూడా వాక్యాలే ఉంటాయి. ఉదాహరణకు,
``ఒకవేళ ఆ సేవకుడు ఆ పని చేసి ఉంటే, తోటమాలి
నిర్దోషి.'' ప్రాథమిక తర్కశాస్త్రంలో
సోపాధికాలను $\lif$ సంయోజకంతో సంకేతీకరించడం
నేర్చుకుంటాం: \dots{} స్థానాల్లోని భాగాలను,
ఉదాహరణకు \tetoken{సూత్రాల}{!!{formula}s}
$!A$,~$!B$తో సూచించి, మొత్తం సోపాధికాన్ని
$!A \lif !B$తో సూచిస్తాం.

$\lif$ సంయోజకం \emph{సత్యమూల్యాధారితం}:
అంటే $!A \lif !B$ యొక్క సత్యమూల్యం---
$\True$ లేదా $\False$---$!A$,~$!B$ల
సత్యమూల్యాల ద్వారానే నిర్ణయించబడుతుంది.
$!A$ అసత్యమైనా, $!B$ సత్యమైనా,
అప్పుడే $!A \lif !B$ సత్యం; మిగిలిన
సందర్భంలో అసత్యం. సత్యమూల్య నియామకం
~$\pAssign v$కు సంబంధించి,
$\pSat/{v}{!A}$ లేదా $\pSat{v}{!B}$
అయినప్పుడే $\pSat{v}{!A \lif !B}$ అని
నిర్వచిస్తాం. ఈ అర్థవిచారంతో ఉన్న
$\lif$ సంయోజకాన్ని \emph{భౌతిక సోపాధికం}
అంటాం.

ఈ నిర్వచనం కొన్ని ప్రాథమిక తార్కిక
వాస్తవాలను ఇస్తుంది. మొదటగా, భౌతిక
సోపాధికానికి నిగమన సిద్ధాంతం వర్తిస్తుంది:
\begin{align}
  \text{ఒకవేళ } \Gamma, !A \Entails !B \text{ అయితే } \Gamma \Entails !A \lif !B
\end{align}
ఇది సత్యమూల్యాధారితం: $!A \lif !B$,
$\lnot !A \lor !B$ తుల్యాలు:
\begin{align}
  !A \lif !B & \Entails \lnot !A \lor !B\\
  \lnot !A \lor !B & \Entails !A \lif !B
  \intertext{భౌతిక సోపాధికం తన ఉత్తరపక్షం
    నుంచీ, పూర్వపక్షం నిరాకరణ నుంచీ
    పర్యవసిస్తుంది:}
  !B & \Entails !A \lif !B\\
  \lnot !A & \Entails !A \lif !B
  \intertext{అసత్య భౌతిక సోపాధికం తన
    పూర్వపక్షం, ఉత్తరపక్షం నిరాకరణల
    సంయోగానికి తుల్యం: $!A \lif !B$
    అసత్యమైతే $!A \land \lnot !B$
    సత్యం; తిరుగుదిశలోనూ అంతే:}
  \lnot(!A \lif !B) & \Entails !A \land \lnot !B\\
  !A \land \lnot !B & \Entails \lnot(!A \lif !B)
  \intertext{భౌతిక సోపాధికం మోడస్
    పోనెన్స్‌కు అనుకూలం:}
  !A, !A \lif !B & \Entails !B
  \intertext{భౌతిక సోపాధికాలను
    ఒకటిగా చేర్చవచ్చు:}
  !A \lif !B,  !A \lif !C & \Entails !A \lif (!B \land !C)
  \intertext{పూర్వపక్షాన్ని ఎల్లప్పుడూ
    బలపరచవచ్చు; అంటే సోపాధికం
    \emph{ఏకదిశ}:}
  !A \lif !B & \Entails (!A \land !C) \lif !B
  \intertext{భౌతిక సోపాధికం సంక్రామకం;
    అంటే గొలుసు నియమం చెల్లుతుంది:}
  !A \lif !B, !B \lif !C & \Entails !A \lif !C
  \intertext{భౌతిక సోపాధికం తన
    ప్రతివిపర్యయానికి తుల్యం:}
  !A \lif !B & \Entails \lnot !B \lif \lnot !A\\
  \lnot !B \lif \lnot !A & \Entails !A \lif !B
\end{align}

ఇవన్నీ గణిత వాదనలో ఉపయోగకరమైన,
సమస్యలేని నిగమనాలు. అయితే గణితేతర
సందర్భాల్లో వీటికి సమానమైన నిగమనాలకు
ఉద్దేశిత ప్రతిదృష్టాంతాలు తాత్త్విక,
భాషాశాస్త్ర సాహిత్యంలో విస్తారంగా
కనిపిస్తాయి. అందువల్ల ఆంగ్లంలోని
``ఒకవేళ \dots అయితే \dots'' వాక్యాలను
సంకేతీకరించేటప్పుడు భౌతిక సోపాధికం
~$\lif$ సరైన సంయోజకం కాదని---లేదా
కనీసం అన్ని సందర్భాల్లో సరైనది కాదని---
సూచిస్తుంది.

\end{document}
